\section{Observabilidade e Debugging}

A arquitetura foi desenhada para permitir localizar uma inconsistência sem depender exclusivamente da inspeção manual do banco de dados.

\subsection{Cadeia de diagnóstico}

Quando um valor exibido no frontend parecer incorreto, a investigação deve seguir a cadeia:

\begin{enumerate}
    \item verificar a resposta da API Java;
    \item comparar com a consulta equivalente no PostgreSQL;
    \item comparar o PostgreSQL com o Gold produzido pelo pipeline;
    \item comparar o Gold com o Silver;
    \item comparar o Silver com o Bronze;
    \item verificar a origem SINAN/PySUS quando necessário.
\end{enumerate}

\placeholderfigure{cadeia-debug}{6cm}{Cadeia de diagnóstico: frontend, Java, PostgreSQL e etapas do pipeline.}

\subsection{API de inspeção}

A API de inspeção do Python pode disponibilizar consultas sobre os artefatos Gold. Ela deve ser tratada como ferramenta operacional e de desenvolvimento, e não como o contrato público da aplicação.

Isso permite uma comparação rápida:

\begin{center}
\texttt{Gold Python} $\leftrightarrow$ \texttt{PostgreSQL} $\leftrightarrow$ \texttt{Java API} $\leftrightarrow$ \texttt{Frontend}
\end{center}

\subsection{Logs}

Os três componentes devem produzir logs suficientes para identificar uma execução sem registrar dados sensíveis desnecessariamente. O pipeline deve registrar, no mínimo, identificadores de execução, origem, período processado, quantidade de registros e resultado das validações.

O Java deve registrar requisições relevantes, falhas de integração e erros de aplicação sem registrar credenciais ou informações sensíveis.

\subsection{Validação automatizada}

Sempre que possível, a comparação entre etapas deve ser automatizada. Exemplos de invariantes:

\begin{itemize}
    \item schema esperado;
    \item quantidade mínima ou esperada de registros;
    \item ausência de chaves geográficas inválidas;
    \item consistência de totais;
    \item equivalência entre dados publicados no PostgreSQL e artefatos Gold.
\end{itemize}

\begin{regra}[Debug por camadas]
A existência de um valor incorreto no frontend não deve levar imediatamente à conclusão de que o frontend está errado. Cada camada deve possuir meios próprios de inspeção para reduzir sistematicamente o espaço de busca do problema.
\end{regra}
